% =====================================================================
%  GMN2 附录 D  多中心 Taub-NUT 上的全纯坐标（Holomorphic coordinates on multi-center Taub-NUT）
%  原文：wkb-paper.tex 9760–9870 行；印刷页 p.172–174；编号公式 (D.1)–(D.15) 共 15 条
%  脚注 0 条；插图 0 幅
%  main.tex 第 67 行 \appendix 生效时，本节自动编号为 D，公式随之为 (D.1)…，与原文一致
%  块内引用（\eqref{eq:holoTN1}）直接用；跨块引用（第 3.2 节、第 3.2.5 节）硬编码原文编号并加 % TODO
%  原文末条 align 的两行 (D.14)(D.15) 按规范拆成两条 equation
% =====================================================================

\section{多中心 Taub-NUT 上的全纯坐标}\label{app:HoloTN}

我们从 Gibbons--Hawking 拟设开始。Taub-NUT 空间 $TN$ 带有一个映射 $\pi: TN \to \IR^3$，其一般纤维是一个圆，度量为
\begin{equation}
ds^2 = V^{-1} \Theta^2 + V (d\vec r)^2,
\end{equation}
其中 $d \Theta = \pi^*( *dV)$。整体良定义的一形式 $\Theta$ 由 $\pi_* \Theta = 4\pi$ 归一。

Taub-NUT 的中心位于 $\vec r_a$，而
\begin{equation}
V = 1 + \sum_a \frac{1}{\vert \vec r - \vec r_a \vert}.
\end{equation}
令 $z=x^1 + i x^2$。某个复结构下的全纯函数被
\begin{equation}
\p_{\bar z} - \Theta_{\bar z} \p_\psi
\end{equation}
\begin{equation}
\p_3 + (i V - \Theta_3) \p_\psi
\end{equation}
所湮灭，其中 $\p_\psi$ 是整体良定义的矢量场，它生成纤维内的旋转，并归一为 $\langle \Theta, \p_\psi \rangle = 1$。

为写出全纯坐标，用片覆盖流形比较方便。对每个中心引入角坐标 $(\theta_a, \phi_a)$。我们将定义片 $\CU_\epsilon$，其中 $\epsilon$ 是一个矢量，对每个中心各有一个分量，且 $\epsilon_a = \pm 1$。我们令 $\CU_\epsilon$ 为使 $\theta_a \not= \half(1-\epsilon_a)\pi$ 的那些 $\vec r$（$\vec r\not=\vec r_a$）的集合。在每个片上，都可以定义周期为 $4\pi$ 的纤维坐标 $\psi_\epsilon$，使得
\begin{equation}
\Theta = d \psi_\epsilon + A_\epsilon
\end{equation}
\begin{equation}
A_\epsilon = \sum_a \epsilon_a (1 + \epsilon_a\cos\theta_a) d \phi_a
\end{equation}
由于 $\Theta$ 整体良定义，由此可以读出 $\psi_\epsilon$ 的过渡函数。现在，定义下面这个量比较方便：
\begin{equation}
R_{a,\pm}: = \vert \vec r - \vec r_a\vert \pm (x_3 - x_{3,a} )
\end{equation}
注意 $R_{a,+}R_{a,-} = \vert z-z_a\vert^2$。直接验证可知
\begin{equation}\label{eq:holoTN1}
\tilde U_\epsilon = \prod_{\epsilon_a=+1} (R_{a,-})^{-1/2}
\prod_{\epsilon_a = -1} (R_{a,+})^{1/2} e^{\half (i \psi_\epsilon + x_3)}
\end{equation}
被反全纯矢量场湮灭。若把 \eqref{eq:holoTN1} 乘上 $\prod_{\epsilon_a=+1}(z-z_a)$，全纯性得以保持；我们正是这样做来定义
\begin{equation}
U_\epsilon = \prod_a (R_{a,+})^{1/2} \prod_{\epsilon_a=+1}
\frac{z-z_a}{\vert z-z_a\vert} e^{\half (i \psi_\epsilon + x_3)}
\end{equation}
可以验证 $U_\epsilon$ 实际上与 $\epsilon$ 无关，也就是说，它是\emph{整体良定义}的。正如我们已经说过的，它在 $TN$ 上是全纯的。它的除子是若干全纯圆盘的不交并，这些圆盘投影到直线 $\theta_a = \pi$ 上。今后我们略去下标 $\epsilon$，直接写作 $U$。$TN$ 上任一绕纤维的绕数为 1 的整体全纯函数，必定是 $z$ 的多项式乘以 $U$。

现在也可以验证
\begin{equation}
W = \prod_a (z-z_a)/U
\end{equation}
是整体良定义且全纯的。它的除子是若干全纯圆盘的并，这些圆盘投影到直线 $\theta_a = 0$ 上。可以显式地写出
\begin{equation}
W = \prod_a (R_{a,-})^{1/2} \prod_{\epsilon_a=-1}\frac{ z-z_a}{\vert z-z_a\vert} e^{-\half ( i \psi_\epsilon + x_3)}.
\end{equation}

在第 3.2 节% TODO 合稿改 \ref{subsec:WittenConstruction}
中，我们考虑固定 $\vec r$、令 $x_{3,a} \to \pm \infty$ 的极限。假设对 $a \in \CR$ 取 $x_{3,a} \to +\infty$，对 $a \in \CL$ 取 $x_{3,a} \to -\infty$。那么 $\vec r$ 落在这样的片里：对 $a \in \CR$ 有 $\epsilon_a = +1$，对 $a\in \CL$ 有 $\epsilon_a = -1$。可以算得
\begin{equation}
\tilde U_\epsilon \to \prod_{a\in \CR} (2 x_{3,a})^{-1/2}
\prod_{a\in \CL} (2 \vert x_{3,a}\vert)^{+1/2} e^{\half (i \psi +x_3)}
\end{equation}
因此，为使极限良好，我们在这个特殊的坐标片里把 $U$ 归一化，使得
\begin{equation}
U = \prod_{a\in \CR} (z-z_a) \prod_{a\in \CR} \left(
\frac{2x_{3,a}}{R_{a,-}} \right)^{1/2} \prod_{a \in \CL}
\left(\frac{ R_{a,+}}{2\vert x_{3,a}\vert}\right)^{1/2} e^{\half (i \psi_\epsilon + x_3)}
\end{equation}

现在，我们可以把这个结果用到第 3.2.5 节% TODO 合稿改 \ref{subsubsec:LinConfQuivMatt}
的讨论中。我们把柱坐标取为 $t = e^{-\half (i \psi_\epsilon + x_3)}$，把 $z$ 与 $v$ 等同，而把 $x^3$ 与 $x^6$ 等同。于是在极限下
\begin{equation}
U \to t^{-1} \prod_{a\in \CR} (v-v_a),
\end{equation}
\begin{equation}
W \to t \prod_{a\in \CL} (v-v_a),
\end{equation}
其中 $(t, v)$ 是 $T^*\IC^\times$ 上的自然全纯坐标。
